\documentclass[11pt]{letter}
\usepackage[utf8]{inputenc}
\usepackage[T1]{fontenc}
\usepackage[margin=1in]{geometry}
\usepackage{hyperref}
\usepackage{siunitx}
\urlstyle{same}

\signature{Jonathan Mace\\Corresponding author}
\address{Jonathan Mace\\Microsoft Research\\Redmond, WA\\jonathanmace@microsoft.com}
\date{\today}

\begin{document}

\begin{letter}{%
Professor A.\ V.\ Metrikine\\
Editor-in-Chief\\
\textit{Journal of Sound and Vibration}}

\opening{Dear Professor Metrikine,}
\small

Please consider our manuscript, ``Scaling Laws for the Flexural Resonance of Fluid-Filled Viscoelastic Shells: Predictions Across Mammalian Scales,'' as a Short Communication for the \textit{Journal of Sound and Vibration}. The paper develops a Buckingham Pi reduction for a fluid-filled viscoelastic shell model and shows that the flexural resonance problem can be expressed in terms of five governing dimensionless groups. A closed-form scaling law is derived from the shell equations and verified against a 1458-point parametric sweep, with collapse to numerical round-off. The result is a compact structural-dynamics framework for comparing resonance behaviour across mammalian scales without re-solving the full eigenproblem for each geometry.

We believe the manuscript is particularly well suited to \textit{JSV} because its central contribution is modal rather than purely biological. The paper identifies the shell quantities that control low-order flexural resonance, clarifies why the dimensionless frequency remains approximately conserved under realistic allometric constraints, and shows that the associated scattering coupling ratio is correspondingly near size-independent. It also demonstrates that the fluid-dominated breathing mode would require an organism of order \SI{20}{\metre} in size to enter the infrasonic band, ruling out that mechanism across biological scales. These results should interest readers in structural dynamics, shell vibration, and modal scaling, while also contributing to our broader research programme on fluid-filled biological structures.

The manuscript is original, has not been published previously, is not under consideration elsewhere, and has been approved by both authors. The co-author contact email for Brian R. Mace is \texttt{b.mace@auckland.ac.nz}. We hope the paper will be of interest to the journal as a concise and analytically grounded contribution to modal scaling in fluid-loaded shells.

\closing{Yours sincerely,}

\end{letter}
\end{document}
